% Generated by benchmarks/build_stochastic_horizon_evidence.py. Do not edit.
\begin{table}[t]
\centering
\caption{Reduced-horizon stochastic SEP evidence}
\label{tab:stochastic-horizon-evidence}
\small
\begin{tabular}{lccc}
\toprule
SEP horizon & Accepted paths & Maximum residual & Newton iterations \\
\midrule
3 & 8/8 & $1.73\times 10^{-12}$ & 5--5 \\
4 & 8/8 & $7.44\times 10^{-9}$ & 4--5 \\
5 & 8/8 & $1.93\times 10^{-10}$ & 5--7 \\
\midrule
Horizon comparison & All-state max & $\log Y$ max & $\pi^{ann}$ / $i^{ann}$ max \\
\midrule
3$\rightarrow$4 & 0.0137 & 0.00029 & 0.000968 / 0.00112 \\
4$\rightarrow$5 & 0.0212 & 0.000396 & 0.00122 / 0.00139 \\
3$\rightarrow$5 & 0.0349 & 0.000686 & 0.00219 / 0.00251 \\
\bottomrule
\end{tabular}

\vspace{0.4em}
\parbox{0.96\textwidth}{\footnotesize\textit{Notes:} Horizons three through five use eight prespecified path evaluations spanning seven unique points in a 12 percent perturbation design for the two policy-rule parameters; the reference calibration is evaluated twice. The calculation uses order-one stochastic extended path, three Gauss--Hermite nodes, and an $L_\infty$ residual gate of $10^{-8}$. The horizon-six hard design fails under both the sequential ($1.02\times 10^{-4}$) and batched ($1.02\times 10^{-4}$) solvers. First-step changes show that residual convergence is not terminal-horizon convergence. This table supports a reduced numerical pilot only; it is not likelihood or posterior evidence.}
\end{table}
